% P16-C3/C4: finite graphical separation and sufficient back-door adjustment.
\section{Graphical separation and adjustment}
\label{sec:p16-graphical-separation-adjustment}

Use the finite structural model and equation-replacement intervention of
\cref{def:p16-finite-scm,def:p16-causal-intervention}. Whenever a
structural model is used below, its disturbances are mutually independent.
The graph includes every endogenous variable whose equation is part of
the model, whether that variable is observed or latent. In particular,
deleting an unobserved common cause from a diagram does not in general
preserve the independent-disturbance factorization used by these results.

The path-blocking definition and sufficient adjustment criterion follow
Pearl's overview \cite[Definition~1, p.~106; Definition~3 and Eq.~(25),
pp.~114--115]{Pearl2009CausalOverview}. The finite proofs, the explicit
ancestral-moral bridge, and the numerical models here are supplied by
the handbook. Identification is a statement about an exact observational
law in a specified model class; estimation from a sample is a further
problem.

\subsection{Paths and their conditioning status}

\begin{hbdefinition}{def:p16-active-path}{Ancestors, paths, and active paths}
In the DAG \(G=(N,E)\), define
\[
 \operatorname{an}_G(J)=J\cup
 \{i:\text{a directed path of positive length leads from }i
                                  \text{ to a member of }J\}.
\]
Thus ancestry includes the starting set, whereas
\(\operatorname{de}_G(J)\) in \cref{def:p16-causal-dag} uses
positive length. A set \(H\subseteq N\) is \emph{ancestral} when
\(\operatorname{an}_G(H)=H\).

An \emph{undirected path in the DAG} is a sequence of distinct vertices
\(i_0,\ldots,i_m\) such that each consecutive pair is joined by an
edge in one of its two directions. The edge directions are retained
when examining the path. An internal vertex \(i_r\), where
\(0<r<m\), is a \emph{collider on this path} precisely when its two
incident path edges have the form
\(i_{r-1}\to i_r\leftarrow i_{r+1}\). Every other internal
vertex is a \emph{noncollider on this path}. These descriptions concern
the two edges used by this path, not all edges incident to the vertex.

For a conditioning set \(C\subseteq N\), the path is
\emph{active given \(C\)} when every internal noncollider is outside
\(C\), and every internal collider belongs to
\(\operatorname{an}_G(C)\). Equivalently, every collider is itself
in \(C\) or has a strict descendant in \(C\). A path that is not
active is \emph{blocked}. For pairwise disjoint \(A,B,C\subseteq N\),
write
\[
 A\mathrel{\perp_G}B\mid C
\]
when there is no active path with one endpoint in \(A\) and the
other in \(B\). This is \emph{d-separation}. Endpoint membership
in \(C\) does not arise under the disjointness hypothesis.

A \emph{walk} permits repeated vertices, but consecutive vertices must
still be adjacent. Collider status is then assigned to each internal
occurrence. The same activity definition applies occurrence by occurrence.
\index{d-separation}\index{collider}\index{ancestor!causal graph}
\end{hbdefinition}

\begin{hblemma}{thm:p16-active-walk-reduction}{proved=annex:p16.active-walk-reduction}{An active walk contains an active path}
For pairwise disjoint \(A,B,C\), an active walk from \(A\) to
\(B\) exists exactly when an active path does. Thus permitting
repeated vertices in the path-blocking convention does not change
d-separation.
\end{hblemma}

\begin{hbdefinition}{def:p16-ancestral-moral-graph}{The ancestral moral graph}
For pairwise disjoint \(A,B,C\), put
\(H=\operatorname{an}_G(A\cup B\cup C)\). Restrict \(G\) to
the vertices \(H\). Its \emph{moral graph}, denoted \(G_H^{\mathrm m}\),
is the undirected graph on \(H\) with an edge between distinct
vertices \(i,j\) if either they are joined by an original directed
edge or they are both parents of some \(k\in H\).
The latter operation joins co-parents. An undirected graph has no
arrow directions and no collider designations.

Delete the vertices \(C\) and their incident edges from this moral
graph. The resulting graph \(G_H^{\mathrm m}\setminus C\)
\emph{separates \(A\) from \(B\)} when no undirected path joins
the two sets. Restricting to ancestors occurs before joining co-parents.
\index{moral graph}\index{ancestral graph}
\end{hbdefinition}

\begin{hbtheorem}{thm:p16-moral-separation-bridge}{proved=annex:p16.moral-separation-bridge}{Path blocking and ancestral moral separation agree}
For every finite DAG and pairwise disjoint vertex sets \(A,B,C\),
\[
 A\mathrel{\perp_G}B\mid C
 \quad\Longleftrightarrow\quad
 G_{\operatorname{an}_G(A\cup B\cup C)}^{\mathrm m}\setminus C
       \text{ separates }A\text{ from }B.
\]
This is a statement about the graph alone.
\end{hbtheorem}

\subsection{The independence supplied by factorization}

The next two results also concern a mass function \(p\) on \(S_N\)
without a specified system of equations. In that setting use the finite
probability space \(S_N\) with law \(p\), and let \(V_i(v)=v_i\)
be its coordinate random variables.

\begin{hbdefinition}{def:p16-graphical-conditional-independence}{Conditional independence of coordinate tuples}
For pairwise disjoint \(A,B,C\subseteq N\), write
\(V_A\mathrel{\perp\!\!\!\perp}V_B\mid V_C\) if, for every
\(c\in S_C\) with \(\mathbb P(V_C=c)>0\), and all
\(a\in S_A,b\in S_B\),
\[
 \mathbb P(V_A=a,V_B=b\mid V_C=c)
 =\mathbb P(V_A=a\mid V_C=c)
       \mathbb P(V_B=b\mid V_C=c).
\]
These are the ordinary finite conditional probabilities. No conditional
value is assigned when \(\mathbb P(V_C=c)=0\). For empty \(C\),
its unique tuple has probability one, giving ordinary independence.
An empty \(A\) or \(B\) also has one tuple and the identity is
automatic.
\end{hbdefinition}

\begin{hblemma}{thm:p16-ancestral-marginal}{proved=annex:p16.ancestral-marginal}{Ancestral marginalization retains the local factors}
Suppose a mass function on \(S_N\) has the product form
\[
 p(v)=\prod_{i\in N}k_i(v_i\mid v_{\operatorname{pa}_G(i)}),
 \qquad \sum_{b\in S_i}k_i(b\mid a)=1
\]
with nonnegative kernels at every parent tuple. For every ancestral
\(H\subseteq N\), its marginal is
\[
 p_H(h)=\prod_{i\in H}k_i(h_i\mid h_{\operatorname{pa}_G(i)}).
\]
The empty marginal and empty product both equal one.
\end{hblemma}

\begin{hbtheorem}{thm:p16-finite-global-markov}{proved=annex:p16.finite-global-markov}{Graphical separation implies conditional independence}
Every mass function factorizing as in
\cref{thm:p16-ancestral-marginal} satisfies, for pairwise disjoint
\(A,B,C\subseteq N\),
\[
 A\mathrel{\perp_G}B\mid C
 \quad\Longrightarrow\quad
 V_A\mathrel{\perp\!\!\!\perp}V_B\mid V_C.
\]
Consequently this implication holds for every finite structural causal
model with mutually independent disturbances, by
\cref{thm:p16-causal-factorization}. The joint law may have zeros.
The conclusion concerns its positive conditioning events.
\end{hbtheorem}

\subsection{A sufficient back-door adjustment theorem}

\begin{hbdefinition}{def:p16-back-door-data}{Treatment, response, and an adjustment set}
Fix one treatment vertex \(t\in N\), a response set
\(B\subseteq N\setminus\{t\}\), and a set
\(C\subseteq N\setminus(\{t\}\cup B)\). Use
\(X=V_t\), \(Y=V_B\), and \(Z=V_C\) for the associated random
variables or tuples. For \(x\in S_t\), the notation \(\{t\}=x\)
for an intervention uses the singleton indexed tuple with coordinate
\(t\) equal to \(x\). Put
\(Y^{t=x}=V_B^{\{t\}=x}\) and
\(Z^{t=x}=V_C^{\{t\}=x}\). When the treatment vertex is named by
its variable \(X\), the superscript \(X=x\) names this same
intervention. A \emph{back-door path} is an undirected path
from \(t\) to a member of \(B\) whose first edge points into
\(t\). The set \(C\) satisfies the \emph{back-door criterion}
for \((t,B)\) when
\[
 C\cap\operatorname{de}_G(\{t\})=\varnothing
 \quad\text{and}\quad
 C\text{ blocks every such back-door path in }G.
\]
For a fixed treatment value \(x\in S_t\), the required
\emph{overlap condition} is
\[
 \mathbb P(X=x,Z=z)>0
 \quad\text{for every }z\in S_C
                   \text{ with }\mathbb P(Z=z)>0.
\]
Only the indicated treatment value is required here. A contrast of two
values requires overlap at both. The graph criterion uses the complete
specified graph, including any latent vertices along its paths.
\index{back-door criterion}\index{causal positivity!overlap}
\end{hbdefinition}

\begin{hbtheorem}{thm:p16-finite-back-door-adjustment}{proved=annex:p16.finite-back-door-adjustment}{Sufficient adjustment with explicit observation and overlap}
Let \(M\) be a finite structural causal model with mutually independent
disturbances. Suppose \(C\) satisfies \cref{def:p16-back-door-data},
and fix \(x\) satisfying its overlap condition. For every
\(z\) with \(\mathbb P(Z=z)>0\), and every \(y\in S_B\),
\[
 \mathbb P(Y^{t=x}=y\mid Z^{t=x}=z)
       =\mathbb P(Y=y\mid X=x,Z=z).
\]
In particular,
\[
 \mathbb P_M(Y=y\mid\operatorname{do}(X=x))
 =\sum_{\substack{z\in S_C\\\mathbb P(Z=z)>0}}
       \mathbb P(Y=y\mid X=x,Z=z)\mathbb P(Z=z).
\]
The treatment is a single vertex with any
finite nonempty value set; the response and adjustment variables may
be tuples.

If \(\{t\}\cup B\cup C\subseteq O\) is observed, this formula
identifies the displayed intervention law from \(p_{M,O}\) within
the class with the specified graph, value sets, independence assumption,
and overlap. Other endogenous variables may remain unobserved. The
claim does not require their observational distribution as an input.
\end{hbtheorem}

The sum omits zero-mass strata instead of multiplying an undefined
conditional probability by zero. The theorem supplies a sufficient
adjustment condition. It does not assert that every identifiable effect
has an adjustment set, or that every admissible adjustment formula must
arise from this criterion.

\subsection{Fully specified numerical models}

\begin{hbdefinition}{def:p16-adjustment-example-data}{A confounded model with a nonzero treatment effect}
Let \(U_Z\) be a fair bit, \(U_X\) a bit with probabilities
\(\mathbb P(U_X=0)=3/4\), \(\mathbb P(U_X=1)=1/4\), and
\(U_Y\) uniform on \(\{1,\ldots,8\}\), all mutually independent.
Concretely the outcome \((z,a,u)\) has mass
\(\frac12\alpha(a)\frac18\), where
\(\alpha(0)=3/4,\alpha(1)=1/4\).
The binary endogenous equations are
\[
 Z=U_Z,\qquad X=Z\oplus U_X,\qquad
 Y=\mathbf1_{\{U_Y\leq1+2X+4Z\}}.
\]
The graph has edges \(Z\to X\), \(Z\to Y\), and \(X\to Y\).
The symbol \(\oplus\) is the bit exclusive-or defined in
\cref{def:p16-confounded-binary-model}. All functions are specified
on their full input domains.
\end{hbdefinition}

\begin{hbtheorem}{thm:p16-adjustment-numerical-example}{proved=annex:p16.adjustment-numerical-example}{Adjustment removes the confounded part of the contrast}
In \cref{def:p16-adjustment-example-data}, the eight joint masses are
\[
\begin{array}{cc|cc}
 Z&X&Y=0&Y=1\\\hline
 0&0&21/64&3/64\\
 0&1&5/64&3/64\\
 1&0&3/64&5/64\\
 1&1&3/64&21/64
\end{array}
\]
The observed conditional response probabilities are
\(\mathbb P(Y=1\mid X=0)=1/4\) and
\(\mathbb P(Y=1\mid X=1)=3/4\). The intervention probabilities are
\[
 \mathbb P(Y^{X=0}=1)=3/8,\qquad
 \mathbb P(Y^{X=1}=1)=5/8.
\]
Thus the observed contrast is \(1/2\), and the average intervention
contrast is \(1/4\). The measured set \(\{Z\}\) satisfies the
back-door criterion and overlap at both values of \(X\); its
adjustment formula gives these same intervention probabilities.
\end{hbtheorem}

\begin{hbdefinition}{def:p16-collider-example-data}{A collider and its observed descendant}
Let \(X=U_X\) and \(Y=U_Y\) be independent fair bits, and set
\[
 K=\mathbf1_{\{X=1\text{ or }Y=1\}},\qquad D=K.
\]
The graph is \(X\to K\leftarrow Y\), \(K\to D\).
The disturbances for \(K,D\) are constants. A concrete outcome
space is \(\{0,1\}^2\) with mass \(1/4\) at each pair;
the two coordinate disturbances and the constant disturbances are
mutually independent.
\end{hbdefinition}

\begin{hbtheorem}{thm:p16-collider-numerical-example}{proved=annex:p16.collider-numerical-example}{Conditioning on a collider can create dependence}
The variables \(X,Y\) in \cref{def:p16-collider-example-data}
are independent. Conditional on either \(K=1\) or \(D=1\), their
joint table is
\[
\begin{array}{c|cc}
 &Y=0&Y=1\\\hline
 X=0&0&1/3\\
 X=1&1/3&1/3
\end{array}
\]
and they are dependent. The path through \(K\) is blocked given
the empty set and active given either \(\{K\}\) or \(\{D\}\).
Under \(\operatorname{do}(K=1)\), however, their joint law remains
the independent uniform table with all four masses \(1/4\).

There is also a model with this same allowed graph and independent
disturbances in which the path is active given \(\{K\}\), yet
\(X\) and \(Y\) are conditionally independent given \(K\).
Hence failure of d-separation is not a dependence theorem for every
model compatible with the graph.
\end{hbtheorem}

\begin{hbmethod}{def:p16-graphical-adjustment-procedure}{Using a graph to justify a conditional or adjustment formula}
Specify the full DAG, the independent-disturbance assumption, the
observed vertices, and the desired treatment and response. For a
conditional-independence query, form the ancestors of its endpoint and
conditioning sets, join co-parents there, delete the conditioning set,
and check undirected separation. The bridge and global Markov theorem
justify the resulting implication. For adjustment, exclude descendants
of the treatment, check all back-door paths, and establish the positive
conditioning events required at the chosen treatment values. Compute
the displayed weighted conditional probabilities only after these
assumptions are specified.
\end{hbmethod}

The front-door construction in \cref{sec:p16-front-door-identification}
uses these separation and adjustment results to identify an effect with
an unobserved common cause. \Cref{sec:p16-do-calculus-soundness} proves
all three do-calculus rules in the finite structural setting.
Completeness of graphical identification and general counterfactual
identification require further constructions. No empirical effect or
sampling guarantee follows from the stipulated numerical models alone.
